% Generated draft: requires mathematical and source review.
\documentclass[11pt]{article}
\usepackage[margin=25mm]{geometry}
\usepackage{fontspec}
\setmainfont{Latin Modern Roman}
\usepackage{amsmath,amssymb,mathtools,unicode-math}
\setmathfont{Latin Modern Math}
\usepackage{enumitem,needspace,xurl,hyperref}
\setlength{\parindent}{0pt}
\setlength{\parskip}{5pt}
\setlength{\emergencystretch}{3em}
\newcommand{\N}{\mathbb{N}}
\newcommand{\Z}{\mathbb{Z}}
\newcommand{\Q}{\mathbb{Q}}
\newcommand{\R}{\mathbb{R}}
\newcommand{\C}{\mathbb{C}}
\newcommand{\F}{\mathbb{F}}
\newcommand{\A}{\mathbb{A}}
\newcommand{\PP}{\mathbb{P}}
\newcommand{\E}{\mathbb{E}}
\newcommand{\eps}{\varepsilon}
\newcommand{\dd}{\,\mathrm{d}}
\begin{document}
\section*{UnsolvedMath ID: 6200040}
\textbf{Source problem code:} catalogue entry
\subsection{Definitions and mathematical statement}
Relative hyperbolicity is data of a group $G$ together with a peripheral family $\mathcal P$. A usual ideal boundary is the Bowditch boundary $\partial(G,\mathcal P)$, equivalently the boundary of an appropriate hyperbolic cusped model under its hypotheses; parabolic points record peripheral subgroups. The peripheral structure and boundary model matter, so an unspecified boundary of the abstract group alone is not selected here.

\textbf{Research program (Uri Bader).} Develop a theory of analysis on ideal boundaries of relatively hyperbolic groups analogous to the theory for hyperbolic groups. This requires identifying the intended metric/conformal data, measure classes, analytic objects, and analogue theorems, while accounting for peripheral/parabolic behavior. The source does not give a fixed theorem list, invariant, function space, or criterion for completion.

In particular, no visual metric, canonical measure, Ahlfors regularity, doubling, or Poincar\'e assumption is supplied for all boundaries. Those would be substantive additional hypotheses rather than automatic parts of the request. Nor is the program replaced by analysis only on a selected $\operatorname{CAT}(-1)$ subclass.

The supplied summary reports developing conformal-boundary analysis by Bourdon--Pajot and others, mentioning $\operatorname{CAT}(-1)$ boundaries, relatively hyperbolic boundaries, and Bourdon's carpet analysis. It gives no exact theorem or reference identifying the scope. The statement is retained as a broad partially developed program, not a single open/solved proposition or a claim that every hyperbolic-group theorem transfers unchanged.
\subsection{Short English statement}
Build boundary analysis for relatively hyperbolic groups, parallel to hyperbolic-group theory and sensitive to peripheral structure. The source does not specify the boundary model, metrics, measures, or target theorems. Reported conformal-analysis progress is partial and uncited in detail.
\subsection{Sources}
\begin{itemize}
\item Local UnsolvedMath registry, permanent ID 6200040, source code catalogue entry, ``Boundaries of Groups and Kleinian Groups --- Problem 40''; proposer: Uri Bader, proposed year 2005. Status partially solved, classification PARTIAL-PROGRESS, updated 2026-08-06.
\item Kapovich, \emph{Problems on Boundaries of Groups and Kleinian Groups} (2005), Problem 40, PDF page 12, \url{https://www.math.ucdavis.edu/~kapovich/EPR/problems.pdf}. Reported PDF-text-normalized extraction and access: 2026-07-29. Current status NEEDS\_REVIEW; the source labels it a problem. Rights: public author-hosted PDF; statement-reuse terms NEEDS\_REVIEW.
\item The summary mentions Bourdon--Pajot and Bourdon's carpet analysis without paper titles, theorem numbers, or precise peripheral/analytic assumptions.
\end{itemize}
\end{document}
